\section{Synchronization in Wireless Sensor Networks: Review}
Paper \cite{djenouri2014synchronization} reviews synchronization mechanisms aimed at Wireless Sensor Networks, in which the single low-cost, low-power nodes should perform synchronization and transmission by themselves. It also underlines the importance of correlation between different reports from different devices to perform data fusion, time-division multiple-access scheduling and to use real-time applications, even if the delay variability of the wireless connection and the computation limitations of the nodes make the issue even more complex. 

\paragraph{General Concepts} The main idea regarding synchronization in this paper is to mathematically estimate the oscillators in each node, and then implement some sort of message exchange mechanism between them. The goal is to model the differences between the clocks in the devices and modify their behaviour accordingly.

The \textbf{clock models} analyzed, on which the synchronization algorithms will base their estimations, are:
\begin{itemize}
    \item Offset-only model. The local clock of a node is defined as in \ref{offset-only} 
    \begin{equation} \label{offset-only}
        C_i = t + \theta_i 
    \end{equation}
    where \(t\) is an ideal clock and \(\theta_i\) the offset relative to that ideal clock.
    \item Joint skew-offset model. This provides long-term estimators, at the cost of augmented calculation complexity. The local clock can be approximated as done in Equation \ref{skew-offset}
    \begin{equation} \label{skew-offset}
        C_i = \alpha_it + \beta_i 
    \end{equation} 
    where \(\alpha\) is the absolute skew and \(\beta\) is the offset of the node clock. 
\end{itemize}

Once the clock model has been defined, the \textbf{message exchange mechanism} should be chosen:
\begin{itemize}
    \item Sender to Receiver: nodes periodically exchange timestamped synchronization messages. This happens in 4 steps. Node \(n_1\) sends a packet with timestamp \(t_1\). The second node \(n_2\) receives it, with timestamp \(t_2\), and replies with timestamp \(t_3\). At the end, \(n_1\) timestamps the reception of the reply with \(t_4\).
    \item Receiver to Receiver: in this case, a reference (defined exploiting the property of the physical broadcast medium) is used to broadcast messages simultaneously to all devices, that will timestamp the reception event as \(t_1\) and \(t_2\). Those samples are the only ones used for synchronization purposes. 
\end{itemize}

Timestamps obtained through these mechanisms are not totally accurate, as the messages go through a path that may vary nondeterministically to different factors: send time (packet generation and transmission time), medium access time (MAC level), propagation time, reception time (packet processing). The receiver to receiver method removes send and access time from this path.

Once the method has been defined, the paper reports three kind of synchronization categories: 1) event ordering, which is the simplest form, 2) synchronization to a reference and 3) relative clocks. The last one is the most used in the analyzed algorithms, and it consists in maintaining inside the node the information about relative skew and offset to other nodes. 

\paragraph{Main issues} The paper proceeds with a list of all the possible problems in maintaining a good time synchronization between the nodes:
\begin{itemize}
    \item Fast drifting. This is considered the main issue, and it derives from the usage of cheap circuitry components (in particular, oscillators). Usually sensor nodes present two kind of oscillator, a high-speed one (usually with high drifting), and a lower speed one, used in energy-saving situations, that is usually more precise but provides less granularity. 

    Of course, a fast drifting clock requires frequent re-synchronizations, that may result very expensive. 
    \item Node limitation. If the synchronization requires some calculation at the node level, usually the calculation power at disposal is very low, and all of it should be optimized as more as possible. For example, some nodes may not dispose of FPU units, and FP calculation may be implemented at software level, causing delays.
    \item Delay Variation. This is intrinsic to the message exchange mechanism, as there is some nondeterminism when considering message delays. Usually, this issue can be overcome by timestamping packets at the MAC layer, but this solution is both not portable between different devices/protocols, and not applicable to devices whose software stack hides the lower level to the developer.
    \item Fault Tolerance. These synchronization mechanism neglect aspects like nodes failure or packet losses, and provide no recovery algorithm in those cases.
    \item Accuracy measurements. The accuracy of such protocols is extremely difficult to assess, as instantaneous local times from the nodes should be captured at the same physical time. 
    \item Impact of empirical parameters. All the possible implementations relay on some parameters regarding the message exchange, in particular, the number of rounds before estimation \(k\) and the period between rounds \(T\). 
    
\end{itemize}

\paragraph{Conclusions}
The main issue found in this paper with regards to time synchronization is, without doubts, the clock drifting issue. Two solutions have been found: frequent message exchange in the simpler offset-only clock model, or less messages using the more computationally expensive skew-offset model. The choice of the final model is up to the developer, that should choose the least expensive of the two considering the node and the connection protocol at their disposal.

The second big challenge is the delay of message exchange/handling. The usage of the Receiver to Receiver model will reduce that delay, but optimizations can be performed using a MAC-layer timestamping solution (when possible).  